\section{Related Work}
\label{sec:related}

\paragraph{Learned reconstruction through scattering media.}
Deep learning entered speckle imaging with the observation that training on a class of diffusers yields networks that generalize statistically across an unseen diffuser of the same class \cite{li2018deep}, replacing deterministic transmission-matrix inversion \cite{popoff2010image}, whose calibration is highly sensitive to speckle decorrelation \cite{li2018deep}. Subsequent work extended the setting to dynamic scattering with real-time inference \cite{liu2024learning}, memory-less architectures for decorrelating media \cite{zhang2024memoryless}, and transfer-learning adaptations across conditions \cite{li2025cascade}. Our work does not propose a new architecture; it uses a minimal residual-U-Net family as an instrument to test a training-data choice that every method in this line shares but none controls.

\paragraph{Fabrication, hallucination, and evaluation in learned imaging.}
That learned reconstructions can invent content has moved from anecdote to mechanism: a recent analysis traces generalization and hallucination in scattering-media networks to physical illumination-zone structure \cite{zhang2026physical}, and a parallel line builds quantitative assessment frameworks for scattering-imaging generalization \cite{long2026universal}. In general generative reconstruction, the Hallucination Index provides a no-reference quality metric reasoning about the forward process \cite{tivnan2024hallucination}. The out-of-band energy ratio used here differs from both: it is anchored to the \emph{measured} MTF support of the system, and it is used as the dependent variable of a causal intervention rather than as a generic quality score. Bandwidth-limited or low-pass filtered training targets are a known practical stabilization trick in deconvolution and restoration; to our knowledge it has never been (i) analyzed as a causal control for spectral fabrication, (ii) quantified against the measured information support of the operator, or (iii) compared against a compute-matched control that receives identical additional optimization---the three elements that separate a folklore practice from a protocol.

\paragraph{Forward-model-informed protocols and linear baselines.}
Classical speckle physics supplies the reference structure our protocol quantifies against: the memory effect bounds shift-invariant reconstruction \cite{freund1988memory}, speckle correlation carries object autocorrelation information \cite{katz2014noninvasive}, and the Wiener filter gives the optimal linear reconstruction under a known transfer function and noise level \cite{wiener1949extrapolation}. Learned methods are frequently compared with linear baselines on fidelity metrics; we additionally compare on the fabrication axis, which exposes that neither dominates---they occupy distinct operating points. Reviews of deep learning in computational imaging have noted qualitatively that super-resolution beyond physical limits depends on learned priors rather than new information \cite{barbastathis2019use}; this paper turns that observation into a measured, protocol-level statement with a controlled cause and a quantitative controller.
